\chapter{Contexte réglementaire et modélisation en assurance vie}
\label{chap:contexte}

\section{Les spécificités des produits d'assurance vie épargne}
\label{sec:spec_av}

\subsection{Principes fondamentaux du contrat d'assurance vie}

L'assurance vie peut désigner plusieurs réalités selon le contexte. Il peut s'agir d'un contrat d'assurance ou d'une des deux grandes branches de l'assurance. De façon générale, elle regroupe les assurances dont le risque est viager, c'est-à-dire lié à la durée de vie de l'assuré. Dans le cadre de ce mémoire, nous définissons l'assurance vie comme un contrat d'assurance dont le risque est porté par la durée de vie de l'assuré et dont le but précis est la constitution d'une épargne. Nous excluons ainsi de notre périmètre les contrats de prévoyance (qui est une forme d'assurance en cas de décès couvrant aussi les risques d'invalidité, d'incapacité et de dépendance). L'assurance vie est donc un produit de long terme, permettant à l'assuré de se constituer un capital ou une rente, tout en bénéficiant, selon les contrats, d'une protection enc as de décès. L'assureur s'engage à verser un capital à l'assuré ou à ses bénéficiaires en cas de décès de l'assuré, ou à verser un capital ou une rente à l'assuré à l'échéance de son contrat. Cette mécanique est la preuve de l'inversion du cycle de production en assurance vie où l'assureur reçoit des primes de l'assuré avant des verser des prestations futures. La vraie incertitude réside au moment du rachat et du montant du rachat. Il est donc primordial pour l'assureur de bien piloter son activité afin d'assurer sa solvabilité et sa liquidité à long terme tout en créant de la valeur. Ce pilotage est l'objet de l'ALM. L'idée est de piloter conjoinctement les actifs et les passifs afin de s'assurer que l'engagement tiendra les engagements futurs.




